\section{A Vandermonde estimate with linear population cost}
\label{sec:vandermonde}

A termwise estimate of pair phases suggests an unnecessary harmonic
sum. Summing the linear part first removes it. This observation is
useful because a cluster has $m(m-1)/2$ pairs, whereas its physical
cap should be compared with a population of size $m$.

\begin{lemma}[Perturbing a clock discriminant]
\label{lem:vandermonde}
Let $c_j=c_1+(j-1)\pi$, $1\le j\le m$, and let
$w_j=c_j+\delta_j$ with $|\delta_j|\le\eta\le\pi/8$.
For $m=1$ both discriminants below are one. Using the logarithm
obtained by continuously moving the points from
$c_j$ to $w_j$, one has
\begin{equation}
 \left|\log\frac{\prod_{i<j}(w_j-w_i)^2}
                      {\prod_{i<j}(c_j-c_i)^2}\right|
 \le \frac{4(m-1)}\pi\eta+\frac{8m}{9}\eta^2
 \le 2m\eta.
 \label{eq:vandermondebound}
\end{equation}
\end{lemma}
\begin{proof}
Set $u_{ij}=(\delta_j-\delta_i)/(\pi(j-i))$, so $|u_{ij}|\le1/4$.
Each pair ratio stays in a disk about one avoiding zero. Hence its
logarithm is unambiguously the power-series logarithm, and
\[
 \log(\Delta(w)/\Delta(c))
 =2\sum_{i<j}u_{ij}+2\sum_{i<j}\{\Log(1+u_{ij})-u_{ij}\}.
\]
With $H_0=0$ and $H_j=\sum_{d=1}^j d^{-1}$,
\[
 \sum_{i<j}\frac{\delta_j-\delta_i}{j-i}
 =\sum_{j=1}^m\delta_j(H_{j-1}-H_{m-j}).
\]
For each $d$, the two indicator sequences in
$H_{j-1}-H_{m-j}$ differ at $2\min(d,m-d)$ indices. Therefore
\[
 \sum_{j=1}^m|H_{j-1}-H_{m-j}|
 \le\sum_{d=1}^{m-1}\frac{2\min(d,m-d)}d
 \le2(m-1).
\]
This bounds the linear part by $4(m-1)\eta/\pi$.
The elementary remainder estimate
$|\Log(1+u)-u|\le |u|^2/[2(1-|u|)]$ bounds the rest by
\[
 \frac{16\eta^2}{3\pi^2}
       \sum_{d=1}^{m-1}\frac{m-d}{d^2}
 \le\frac{8m}{9}\eta^2.
\]
Finally $4/\pi+\pi/9<2$, proving the last inequality.
\end{proof}

The estimate controls a continuously lifted complex logarithm, not
only the modulus of a determinant. A lower gap bound by itself
would not give the same phase information. The proof is elementary;
no optimality or priority for its numerical constants is asserted.

\section{The complete theta source on a narrow spatial strip}
\label{sec:source}

We now prove Theorem~\ref{thm:headline}. The central simplification
is to work in disks around the clock points, whose imaginary parts
are bounded independently of the rank. On such disks the exponential
bound for sine and cosine involves $|\im w|$, not $|w|$. The real
width can therefore grow without introducing an $e^{\pi K}$ loss.

\begin{lemma}[Weighted concentration]
\label{lem:moments}
Let $U$ be the positive random variable with density
$e^{Tu^2}\Phi(u)/H(T,0)$ on $u\ge0$, and put $V=U-L$.
For each fixed $a\ge0$ and $j=1,2,3$,
\begin{equation}
 M_j(T,a):=\mathbb E\bigl[|V|^j e^{a|V|/L}\bigr]
 =O_a(\kappa^{-j/2}).
 \label{eq:moments}
\end{equation}
The constants do not depend on a root rank or a selected population.
\end{lemma}
\begin{proof}
All theta indices can be retained in the factorization
\begin{align*}
 e^{Tu^2}\Phi(u)&=2\pi^2 e^{q_T(u)}b(u),&
 q_T(u)&=Tu^2+9u-\pi e^{4u},\\
 b(u)&=\sum_{n\ge1}\left(n^4-\frac{3n^2e^{-4u}}{2\pi}\right)
          e^{-\pi(n^2-1)e^{4u}}.
\end{align*}
For all $u\ge0$,
\begin{equation}
 0<b_*:=1-\frac3{2\pi}\le b(u)
 \le B_*:=\frac1{1-16e^{-3\pi}}.
 \label{eq:sourcebounds}
\end{equation}
The lower bound uses the first term; all terms are positive. For
the upper bound discard the subtraction and compare successive terms
$n^4e^{-\pi(n^2-1)}$ with a geometric series of ratio at most
$16e^{-3\pi}$.

For large $T$, $L\to\infty$, $T/\kappa\to0$, and
\[
 q_T(L+v)-q_T(L)=Tv^2-\pi e^{4L}(e^{4v}-1-4v).
\]
On $v\ge-1/4$ one has $q_T''(L+v)\le-\kappa/(2e)$
eventually. Hence this difference is at most
$-c\kappa v^2$, where $c=1/(4e)$.
On $0\le u\le L-1/4$, $q_T$ has at most one local minimum and
no local maximum: $q_T'$ first increases and then decreases, with
its later zero at $L$. Since $q_T(L-1/4)>q_T(0)$ eventually,
the maximum on that interval is its right endpoint. Thus the entire
far-left part has exponent at most $q_T(L)-c\kappa/16$.

On $|v|\le\kappa^{-1/2}$, Taylor's formula gives
$q_T(L+v)-q_T(L)\ge-1/2-O(\kappa^{-1/2})$.
The normalization is therefore at least a positive constant times
$e^{q_T(L)}/\sqrt\kappa$. On the near and right parts put
$w=\sqrt\kappa v$. Dividing by this lower normalization bounds
the moment by a constant times
\[
 \kappa^{-j/2}\int_\R |w|^j
 e^{-cw^2+a|w|/(L\sqrt\kappa)}\dd w
 +L^{j+1}\sqrt\kappa\,e^{a-c\kappa/16}.
\]
The integral is bounded for large $T$; the second term is smaller
than every negative power of $\kappa$. This proves the assertion.
Both tails, the normalization, and all source indices were included.
\end{proof}

\begin{lemma}[A rank-uniform complex-time approximation]
\label{lem:strip}
For real $|s|\le Y$ define
\[
 F(s,w)=\frac{H(T+is,w/L)}{H(T,0)},\qquad
 c_j(w)=\partial_w^j\cos w.
\]
On $|\im w|\le a$ and $|w|\le Z$, for $j=0,1,2,3$,
\begin{equation}
 |\partial_w^jF-e^{isL^2}c_j(w)|\le
 e^a\left\{
 \sum_{k=1}^j\binom jk\frac{M_k(T,a)}{L^k}
 +\frac{ZM_1(T,a)}L
 +Y\bigl(2LM_1(T,a)+M_2(T,a)\bigr)\right\}.
 \label{eq:stripbound}
\end{equation}
The sum is zero for $j=0$. In particular, for fixed $a,Y$ and
$Z\le\pi K+1$, the right side is $O_{a,Y}(\varepsilon_T)$.
\end{lemma}
\begin{proof}
Put $R=U/L$. Differentiation under the convergent complete source
gives $\partial_w^jF=\mathbb E[R^je^{isU^2}c_j(Rw)]$.
Split the error into the changes in $R^j$, in $c_j(Rw)$, and in
$e^{isU^2}$. The first is bounded using
\[
 |R^j-1|\le\sum_{k=1}^j\binom jk|V|^k/L^k,
 \qquad |c_j(Rw)|\le e^a e^{a|V|/L}.
\]
The line segment from $w$ to $Rw$ gives
$|c_j(Rw)-c_j(w)|\le Z|V|e^ae^{a|V|/L}/L$.
Finally, because $s,U,L$ are real,
\[
 |e^{isU^2}-e^{isL^2}|
 \le |s|\,|U^2-L^2|\le Y(2L|V|+V^2).
\]
Taking expectations proves \eqref{eq:stripbound}. No complex-time
derivative of an asymptotic error has been taken.
\end{proof}

\subsection{Actual moving disks and their rank labels}

Let $s_j=(j-1/2)\pi$, $r_1=1/3$, and $\rho=1/4$.
For $|v|\le r_1$ consider
\begin{equation}
 F(s,w)-vF_w(s,w)=0,
 \qquad w_{0,j}(v)=s_j+\arctan v.
 \label{eq:rootinverse}
\end{equation}
The principal branch of $\arctan$ is analytic here. On
$|w-w_{0,j}(v)|=\rho$, the reference equation is
\[
 e^{isL^2}(\cos w+v\sin w)
 =-\sin(s_j)e^{isL^2}\sqrt{1+v^2}\,
       \sin(w-w_{0,j}(v)).
\]
Its modulus is at least
$\sqrt{1-r_1^2}(2\rho-\sinh\rho)>0$, uniformly in $j$.
These disks have $|\im w|\le\operatorname{arctanh}(1/3)+1/4<1$
and $|w|\le\pi K+1$ for $j\le K$.
Lemma~\ref{lem:strip} and Rouch\'e's theorem therefore give exactly
one simple solution in each disk, once $\varepsilon_T$ is small.
Throughout the disk,
\[
 |\sin w|=|\cos(w-s_j)|
 \ge2-\cosh\bigl(\operatorname{arctanh}(1/3)+1/4\bigr)>0.
\]
Thus $F_w\ne0$ at the actual solutions as well.
The solutions depend holomorphically on $\tau=T+is$ and $v$
locally, and uniqueness patches them along the compact strip.
Since $H/H_z=v/L$, they are the required inverse disks in the
physical coordinate $q=v/L$.
Injectivity in $v$ follows by recovering $v$ from that quotient.
Their images are disjoint, since the disks around adjacent clock
centers have total diameter less than $\pi$.

At $s=0$ the source and the reference are real. Choose any fixed
$0<\beta<\pi/4$. On the real interval from zero to $s_K+\beta$,
the uniform errors for $F$ and $F_w$ tend to zero. On each
$[s_j-\beta,s_j+\beta]$ the derivative has constant nonzero
sign and the two endpoint values have opposite signs. On the
remaining intervals $|\cos w|\ge\sin\beta$, so no zero is
missed. This identifies the roots as exactly the first $K$
positive real roots, not merely $K$ roots near guessed labels.

At the centers $v=0$, the same estimates imply, uniformly for
$j\le K$ and $|s|\le Y$,
\begin{align}
 Lz_j(T+is)&=s_j+\delta_j(s),& |\delta_j(s)|&\le C_Y\varepsilon_T,
 \label{eq:rooterror}\\
 h_j(T+is)&=H(T,0)L[-\sin(s_j)]e^{isL^2}(1+\rho_j(s)),&
 |\rho_j(s)|&\le C_Y\varepsilon_T.
 \label{eq:herror}
\end{align}
For \eqref{eq:rooterror}, use the lower linear bound for $\sin$
on the fixed root disk. For \eqref{eq:herror}, evaluate the
$F_w$ estimate there and use the bounded derivative of $\sin$.
Increasing $C_Y$ if needed pays both bounds; it is independent of
the population. The ratio in \eqref{eq:herror} remains in a disk
about one of radius less than $1/2$, so it has a single near-one
logarithm throughout the cap. This is what controls a lifted phase,
even though $YL^2$ itself becomes large.

\subsection{Proof of the terminal estimates}

Equation~\eqref{eq:herror}, summed over the selected block, gives
\begin{equation}
 \Theta(T,Y)=mYL^2+O_Y(m\varepsilon_T).
 \label{eq:thetaerror}
\end{equation}
Substitution in \eqref{eq:closures}, with
$k_\epsilon=r^2/(L^2+r^2)$, proves \eqref{eq:capheadline}.
This argument uses endpoint logarithms and the exact disk integral;
it does not multiply an uncontrolled derivative error by a growing
cap volume.

Apply Lemma~\ref{lem:vandermonde} to the unfolded roots in
\eqref{eq:rooterror}. The reference discriminant is positive real
and independent of $s$; the factor $L^{-m(m-1)}$ is also positive
real. For $\eta=C_Y\varepsilon_T\le\pi/8$ the two endpoint
logarithms have modulus at most $2m\eta$ each. Therefore
\begin{equation}
 \left|\frac{k_\epsilon}{2\pi}
       \im\int_{T-iY}^{T+iY}d\log\Delta\right|
 \le\frac{2k_\epsilon m\eta}{\pi}.
 \label{eq:vfinite}
\end{equation}
This proves \eqref{eq:vheadline}. The factorization
\eqref{eq:factor} proves the corresponding analytic-factor estimate
by subtraction. In particular,
\begin{equation}
 \frac{k_\epsilon}{2\pi m}
       \im\int_{T-iY}^{T+iY}d\log\UU
 =\frac{r^2YL^2}{\pi(L^2+r^2)}
       +O_Y\!\left(\frac{r^2\varepsilon_T}{L^2}\right).
 \label{eq:unitbound}
\end{equation}

If $Y=Y_T=S/L^2$ with fixed $S>0$, repeat
\eqref{eq:stripbound} with that value, rather than using a fixed-width
error. Its scale becomes $O_S((K+1)/(L\sqrt\kappa))$.
The same construction gives
\begin{equation}
 \frac{\II_{T,m}}m
 =\frac{r^2S}{\pi(L^2+r^2)}
 +O_S\!\left(\frac{r^2(K+1)}{L^3\sqrt\kappa}\right)
 =O_{r,S}(L^{-2}).
 \label{eq:shrinkingcap}
\end{equation}
This is a uniform estimate per root even when $m$ grows. An absolute
vanishing cap follows, for example, if also $m=o(L^2)$.

\subsection{Proof of the late edge estimate}

On $[\theta T,T]$, $L_t/L_T\to1$ and
$\kappa_t/\kappa_T$ stays bounded above and below by positive
constants depending on $\theta$. The preceding error estimates
are therefore uniform there. For the fixed physical radius
$\epsilon=r/L_T$, the local inverse variable is
$v=L_tq$, with $|v|\le r$. The same construction gives disjoint
holomorphic tubes over the whole rectangle: the unique inverse
branches agree on overlaps. No assertion about continuation back
to time zero is needed.

At each endpoint $b\in\{\theta T,T\}$, formula
\eqref{eq:thetaerror} reads
$\II_b/m=k_\epsilon YL_b^2/\pi+O_{Y,\theta}(k_\epsilon\varepsilon_T)$.
Thus \eqref{eq:closures} yields the more precise finite asymptotic
\begin{equation}
 \frac{\JJ_\pm}{m}
 =\frac{r^2Y(L_{\theta T}^2-L_T^2)}{2\pi(L_T^2+r^2)}
 +O_{Y,\theta}\!\left(\frac{r^2\varepsilon_T}{L_T^2}\right).
 \label{eq:lateedge}
\end{equation}
Differentiating the defining saddle equation gives the exact identity
$L_t'=2L_t/\kappa_t\sim1/(4t)$, uniformly on this interval.
Hence $L_T-L_{\theta T}=\frac14\log(1/\theta)+o(1)$.
Substitution proves \eqref{eq:edgeheadline} and completes
Theorem~\ref{thm:headline}.

These tubes contain no collision. Their late edge estimate is a
proved bound on a named part of a longer boundary, not an estimate
on the earlier connection to an event-containing region. The latter
has a different asymptotic behavior, described next.
